\section{Conclusion}

This thesis treats deployment surface as a first-class condition in LLM audits.
The core argument is simple: bare API access is not automatically a proxy for the
official consumer-app experience. APIs are reproducible and essential for
controlled evaluation, but consumer applications can expose related model
families through additional product layers, visible sources, local context,
formatting behavior, tools, retrieval, and different refusal styles. If those
layers change the observable output, then the deployment surface must be part of
the audit design.

The methodological contribution is a reproducible protocol for comparing API and
consumer-app behavior using fixed prompts, structured logs, screenshot-backed app
evidence, transcript QA, visible UI-surface coding, pairwise manual review, and
conservative statistical summaries. The thesis also contributes a practical
coding variable, \texttt{ui\_surface\_signature}, for recording visible
app-specific affordances without inferring hidden mechanisms.

The empirical contribution is a completed 60-prompt bare OpenAI API versus
productized ChatGPT app comparison with full pairwise coverage. Manual review of
corrected prompt-echo-cleaned transcripts finds 17 substantive observable
divergences, but the thesis does not treat all 17 as equally strong. The core
behavioral subset is smaller: six refusal-boundary cases, three factual-support
cases, and one format-compliance case. The strongest corrected patterns are:

\begin{itemize}
  \item on refusal-expected prompts, the API more often gives bare full refusals,
  while the app gives bounded safe redirects;
  \item the app exposes visible sources and account-localized context absent from
  plain API text responses, which are retained as product-surface findings rather
  than the main behavioral result;
  \item on short factual controls, the app's retrieval-backed or source-visible
  answers are sometimes coded as supported where the API's memory-based answer is
  coded as a possible mismatch;
  \item after correction, socially sensitive BBQ-style rows do not support an
  API-over-refusal claim and instead serve as negative controls.
\end{itemize}

These findings support the thesis claim that deployment surface is
audit-relevant. They do not support stronger causal claims about access channel
alone, hidden prompts, tools, retrieval, internal guardrails, or routing rules.
The primary statistical endpoint is directional but not conventionally
significant, so the thesis characterizes bounded case evidence rather than
claiming a confirmed population-level effect.

The broader lesson is methodological. LLM evaluations should state the deployment
surface they measure and avoid generalizing API-only findings to consumer-app
settings without validation. This is especially important for safety,
compliance, factuality, and fairness work, where the user-facing answer matters
as much as the underlying model family. A benchmark score collected through an
API may be valid for the API condition and still incomplete as evidence about
ordinary user experience.

Auditing research often faces a choice between scale and evidential care. This
thesis deliberately chose care: a frozen diagnostic bank, screenshot-preserved
app evidence, a correction audit that reduced its own headline count, repetition
checks at both the label and text level, sensitivity analyses, and explicit
claim boundaries. The result is a small but defensible body of evidence for a
simple proposition with practical consequences: the deployment surface through
which a language model is reached is part of the system under audit. Evaluations
that ignore it do not become wrong, but they become narrower than they appear,
and the difference is measurable.
